\subsubsection{Rabin-Karp (single hash)}

\paragraph{Idea intuitiva.}
El algoritmo de \textbf{Rabin-Karp} busca ocurrencias de un patron $P$ de longitud $m$ dentro de un texto $T$ de longitud $n$ deslizando una ventana de tamano fijo $m$ y manteniendo el valor de su hash polinomial. 

Al avanzar la ventana un caracter hacia la derecha, el nuevo hash se calcula en $O(1)$ removiendo el caracter saliente $T[i]$ (multiplicado por $\text{BASE}^{m-1}$) y anadiendo el nuevo caracter entrante $T[i+m]$:
\[
H_{\text{sig}} = \left( (H_{\text{act}} - T[i] \cdot \text{BASE}^{m-1}) \cdot \text{BASE} + T[i+m] \right) \pmod M
\]
Cuando el hash de la ventana coincide con el hash del patron ($tHash == pHash$), se realiza una verificacion caracter por caracter para garantizar que no se trata de una colision. Dado un modulo primo $M \sim 10^9$, la probabilidad de colision aleatoria es insignificante ($1/M$), alcanzando una complejidad lineal esperada $O(n + m)$.

\paragraph{Propiedades clave (invariantes).}
\begin{itemize}\setlength\itemsep{0pt}
	\item Si los hashes son distintos, las cadenas son garantizadamente distintas.
	\item Si los hashes coinciden, se verifica la igualdad exacta para evitar falsos positivos.
	\item Factor multiplicativo $h = \text{BASE}^{m-1} \pmod M$: pondera el caracter mas significativo de la ventana de longitud $m$.
	\item Desplazamiento en $O(1)$: mantener la ventana no requiere rehacer la suma polinomial completa.
\end{itemize}

\paragraph{Codigo clave.}
Desplazamiento del hash de la ventana y verificacion:
\begin{verbatim}
for (int i = 0; i <= n - m; ++i) {
    if (pHash == tHash) {
        bool match = true;
        for (int j = 0; j < m; ++j)
            if (text[i + j] != pattern[j]) { match = false; break; }
        if (match) occ.push_back(i);
    }
    if (i < n - m) {
        tHash = ((tHash - text[i] * h % MOD + MOD) * BASE + text[i + m]) % MOD;
    }
}
\end{verbatim}

\paragraph{Complejidad.}
\begin{itemize}\setlength\itemsep{0pt}
	\item \textbf{Tiempo esperado}: $O(n + m)$ (la verificacion $O(m)$ se ejecuta casi exclusivamente cuando hay una coincidencia real).
	\item \textbf{Peor caso}: $O(n \cdot m)$ si los datos de entrada generan colisiones constantes (por ejemplo cadenas con caracteres todos iguales bajo un modulo vulnerable).
	\item \textbf{Memoria}: $O(1)$ adicional (mas el vector de posiciones de salida).
\end{itemize}

\paragraph{Donde tocar para modificar.}
\begin{itemize}\setlength\itemsep{0pt}
	\item \textbf{Doble hash o base aleatoria}: para evitar casos maliciosos (*anti-hash test cases*), generar \verb|BASE| aleatoriamente o usar dos modulos independientes sin verificacion manual posterior.
	\item \textbf{Busqueda de multiples patrones de igual longitud}: almacenar los hashes de los patrones en un \verb|unordered_set<long long>| para buscarlos todos simultaneamente en una sola pasada $O(n)$.
	\item \textbf{Matching en 2D (matrices)}: calcular hashes por filas y luego por columnas con ventanas de tamano fijo $R \times C$.
\end{itemize}

\paragraph{Trampas y errores frecuentes.}
\begin{itemize}\setlength\itemsep{0pt}
	\item \textbf{Negativos en el operador \texttt{\%}}: al restar \verb|text[i] * h|, el valor puede ser negativo. Se debe sumar \verb|+ MOD| antes de multiplicar por \verb|BASE| y aplicar modulo.
	\item \textbf{Desbordamiento de 64 bits}: la expresion \verb|(tHash - ...) * BASE| puede desbordar si no se aplica modulo intermedio. Usar \verb|long long| para todas las variables involucradas.
	\item \textbf{Caso $|P| > |T|$}: si $m > n$, el bucle no debe ejecutarse; se debe retornar vacio inmediatamente al inicio de la funcion.
\end{itemize}

\paragraph{Explicacion linea por linea.}
\textbf{Funcion rabinKarp:}
\begin{itemize}\setlength\itemsep{0pt}
	\item \verb|const long long BASE = 91138233LL, MOD = 1000000007LL;| -- base prima elegida para dispersion uniforme y modulo primo grande de 32 bits.
	\item \verb|if(m > n) return {};| -- descarta trivialmente cuando el patron es mas largo que el texto.
	\item \verb|for(int i = 0; i < m - 1; i++) h = h * BASE % MOD;| -- calcula $h = \text{BASE}^{m-1} \pmod{\text{MOD}}$, el coeficiente del caracter inicial de la ventana.
	\item \verb|pHash = (pHash * BASE + pattern[i]) % MOD;| -- evalua el hash del patron completo.
	\item \verb|tHash = (tHash * BASE + text[i]) % MOD;| -- evalua el hash de la primera ventana de longitud $m$ en el texto.
	\item \verb|if(pHash == tHash)| -- detecta posible coincidencia entre la ventana actual y el patron.
	\item \verb|for(int j = 0; j < m; j++) if(text[i+j] != pattern[j]) ...| -- verificacion caracter a caracter para garantizar ausencia de falsos positivos por colision.
	\item \verb|if(match) occ.push_back(i);| -- registra el indice de inicio de la ocurrencia confirmada.
	\item \verb|tHash = ((tHash - text[i] * h % MOD + MOD) * BASE + text[i+m]) % MOD;| -- actualiza el hash de la ventana en $O(1)$: elimina $T[i]$ a la izquierda y anade $T[i+m]$ a la derecha.
\end{itemize}
\textbf{Programa principal:}
\begin{itemize}\setlength\itemsep{0pt}
	\item \verb|auto occ = rabinKarp("hello world", "world");| -- busca el patron ``world'' dentro de ``hello world''.
	\item \verb|cout << "found at " << pos;| -- imprime la posicion encontrada (indice 6).
\end{itemize}

\paragraph{Codigo.}
Ver \texttt{cpp.tex}, seccion \textit{Rabin-Karp (single hash)}.

\vspace{0.5em}
